\documentclass[11pt]{article}
\usepackage[margin=1in]{geometry}
\usepackage[T1]{fontenc}
\usepackage{lmodern,amsmath,amssymb,amsthm,booktabs,microtype}
\usepackage[hidelinks]{hyperref}
\newtheorem{theorem}{Theorem}
\newtheorem{lemma}[theorem]{Lemma}
\newtheorem{corollary}[theorem]{Corollary}
\newcommand{\Q}{\mathbb Q}
\newcommand{\C}{\mathbb C}
\newcommand{\Z}{\mathbb Z}
\newcommand{\N}{\mathbb N_0}
\newcommand{\gen}{\operatorname{Gen}}
\newcommand{\Res}{\operatorname*{Res}}
\newcommand{\calM}{\mathcal M}
\setlength{\parindent}{1em}
\setlength{\emergencystretch}{1em}
\title{How many arithmetic progressions resolve a Dirichlet zero?}
\author{}
\date{26 November 2026}
\begin{document}
\maketitle
\vspace{-1.6em}
\begin{abstract}
At a zero shared by Dirichlet $L$-functions, the residues of prime-counting Dirichlet series encode the individual orders of vanishing. We determine the least number of ordinary residue classes needed to recover an arbitrary multiplicity vector, and the sharp precision required for every choice of classes. The precision exponent is a maximum over cyclic subgroups. Each additional inverse pair of generators supplies another pair of embeddings of the corresponding cyclotomic field. Modulo $11$, four classes suffice for exact data, five leave the worst-case precision scale unchanged, and six give uniform separation.
\end{abstract}

\section{What the progressions measure}
Fix $q\ge2$, put $G=(\Z/q\Z)^\times$, $h=|G|$, and write $\widehat G$ for its character group. For a reduced residue $a$ define, initially in $\Re s>1$,
\[
 D_a(s)=\sum_{n\equiv a\bmod q}\frac{\Lambda(n)}{n^s}
 =-\frac1h\sum_{\chi\in\widehat G}\overline{\chi(a)}
                       \frac{L'(s,\chi)}{L(s,\chi)} .
\]
The sum is over positive integers. The identity follows by character orthogonality and the Euler product, and its right side gives the meromorphic continuation. For a specified point $\rho$ in $0<\Re s<1$, put $m_\chi=\operatorname{ord}_\rho L(s,\chi)\ge0$, including zero if the function is nonzero there. Then
\begin{equation}
 -h\Res_{s=\rho}D_a(s)=F_m(a),\qquad
 F_m(a)=\sum_{\chi\in\widehat G}m_\chi\overline{\chi(a)}. \label{eq:F}
\end{equation}
Cancellation can make a measured residue zero. The identity applies to imprimitive characters since the extra Euler factors have no zeros in the open strip. Dirichlet GRH places every nonzero multiplicity vector on $\Re\rho=1/2$. The inverse problem in \eqref{eq:F} also makes sense without a zero-location theorem.

Kisilevsky and Rubinstein explicitly group common zeros with multiplicities in the forward formula for prime races~\cite[\S2.3]{KR}. We ask for the individual multiplicities from selected progression residues. Our observations are complex residues at a known point. Estimating these residues from a finite list of primes would require an additional analytic error bound.

\newpage
\section{A criterion for exact recovery}
The following argument works for every finite abelian group $G$. For a cyclic subgroup $H\le G$, write $d_H=|H|$, $D_H=\varphi(d_H)$ and let $\gen H$ be its generators, including the identity for the trivial subgroup. Fix one generator $a_H$ for each $H$. Write $K_d=\Q(e^{2\pi i/d})$.

\begin{lemma}\label{lem:blocks}
The map from $\Q^{\widehat G}$ sending $v$ to $(F_v(a_H))_H$ is a rational linear isomorphism onto $\prod_H K_{d_H}$. The values at the generators of $H$ are all the conjugates of $F_v(a_H)$, indexed by the embeddings of $K_{d_H}$.
\end{lemma}
\begin{proof}
Each $\chi(a_H)$ is a $d_H$-th root of unity. For $(k,d_H)=1$, applying the substitution $e^{2\pi i/d_H}\mapsto e^{2\pi i k/d_H}$ gives $F_v(a_H^k)$ since $v$ is rational. If all the displayed coordinates vanish then $F_v$ vanishes on every element of $G$, giving $v=0$ by Fourier inversion. Finally the sizes of the disjoint generator sets sum to $|G|$. Hence the target has rational dimension $\sum_H D_H=|G|=|\widehat G|$, equal to the source dimension.
\end{proof}

For a set $A\subseteq G$ of observed classes put $F_v|_A=(F_v(a))_{a\in A}$. By the lemma,
\begin{equation}
 \dim_{\Q}\ker(v\mapsto F_v|_A)
    =\sum_{H:\,A\cap\gen H=\varnothing}D_H. \label{eq:kernel}
\end{equation}
Thus $A$ recovers every nonnegative integer vector from exact values if and only if it meets every generator set. Indeed, any nonzero rational kernel vector can be multiplied by a common denominator and written as a difference of two distinct nonnegative integer vectors. The minimum number of ordinary progression classes is the number of cyclic subgroups. Lemma~\ref{lem:blocks} is the Fourier description of the classical rational abelian group algebra~\cite{PW,CO}.

\section{Sharp precision for a fixed sample}
For integer $M\ge1$, write $\calM_M=\{m\in\N^{\widehat G}:\sum_\chi m_\chi\le M\}$. Measure the resolution of a set $A$ by its minimum spacing
\[
 \delta_A(M)=\min_{m\ne n\in\calM_M}\ \max_{a\in A}|F_m(a)-F_n(a)|.
\]
For an informative $A$ (one meeting every $\gen H$), let $b_H$ count the distinct pairs $\{a,a^{-1}\}$ hit by $A$ inside $\gen H$, for $d_H>2$. Set
\[
 \kappa(A)=\max\!\left(\{0\}\cup
       \left\{\frac{D_H}{2b_H}-1\;:\; d_H>2\right\}\right).
\]
An inverse pair contributes one independent complex observation: $F_m(a^{-1})=\overline{F_m(a)}$.

\newpage
\begin{theorem}[Precision exponent]\label{th:sep}
For every informative $A$ and integer $M\ge1$,
\[
 (2M)^{-\kappa(A)}\ \le\ \delta_A(M)\ \le\ C_{G,A}M^{-\kappa(A)} .
\]
If $A$ misses a generator set, $\delta_A(M)=0$ for all sufficiently large $M$.
\end{theorem}
\begin{proof}
Take distinct $m,n\in\calM_M$ and set $v=m-n$. At least one block $\alpha=F_v(a_H)$ is nonzero. It is an algebraic integer. Every conjugate has modulus at most $2M$ by the triangle inequality. Write $d=d_H$, $D=D_H$, and $\epsilon=\max_{a\in A}|F_v(a)|$. If $d\le2$, the block is an observed nonzero rational integer, so $\epsilon\ge1$. For $d>2$, at least $2b_H$ of the $D$ embeddings give values of modulus at most $\epsilon$. The nonzero algebraic norm is a rational integer, giving
\[
  1\le |N_{K_d/\Q}(\alpha)|\le\epsilon^{2b_H}(2M)^{D-2b_H}.
\]
This proves the lower bound, also when different embeddings give identical values.

If $\kappa(A)=0$, a unit vector and zero show $\delta_A(M)\le1$. Otherwise choose a block attaining $\kappa(A)>0$, and abbreviate $b=b_H$. By Lemma~\ref{lem:blocks}, the subspace supported on this single block has rational dimension $D$. Clear denominators in a basis of its inverse image to get independent vectors $v_1,\ldots,v_D\in\Z^{\widehat G}$. Write $B=\sum_j\|v_j\|_1$. The $(L+1)^D$ vectors $u=\sum_j n_jv_j$, $0\le n_j\le L$, give distinct values on the $b$ observed inverse pairs, since evaluation at even one of the embeddings of a field is injective. Choose one observation per pair. Their real and imaginary parts lie in $[-BL,BL]^{2b}$.

Partition the box into $J^{2b}$ equal cubes, with $J=\lceil(L+1)^{D/(2b)}\rceil-1$. For $L\ge1$, two images lie in one cube. Their difference $w$ satisfies $\|w\|_1\le BL$ and its observed complex coordinates have modulus $O_{G,A}(L^{1-D/(2b)})$. All other blocks vanish. Writing $w=w^+-w^-$ and taking $L=\lfloor M/B\rfloor$ provides the required pair for large $M$. A larger constant covers the finitely many smaller $M$. The last assertion follows from \eqref{eq:kernel}.
\end{proof}

\begin{corollary}[Optimal number of classes]
For $\theta\ge0$, the smallest $|A|$ such that $\delta_A(M)\gg_{G,A}M^{-\theta}$ as $M\to\infty$ equals
\[
 \#\{H:d_H\le2\}
   +\sum_{H:\,d_H>2}\left\lceil\frac{\varphi(d_H)}{2(1+\theta)}\right\rceil .
\]
\end{corollary}
\begin{proof}
Theorem~\ref{th:sep} requires exactly $b_H\ge D_H/[2(1+\theta)]$ in each nonreal block, and one sample in each real block. The independent choices of generators attain their sum.
\end{proof}

\smallskip
\noindent\textbf{Explicit near-collisions.} Modulo $11$ take the block $H=\langle4\rangle$ of order five. Put $\zeta=e^{2\pi i/5}$, $\lambda=(1+\sqrt5)/2$ and $\alpha=(\zeta+\zeta^{-1})^k=\lambda^{-k}$, $k\ge1$. Use $\alpha$ in that block, indexed by $4$, and zero elsewhere. Multiply its inverse Fourier vector by $10$ to obtain $v$. Each entry is a sum of conjugates of algebraic integers, rational by Lemma~\ref{lem:blocks}, hence integral. The other real conjugate of $\zeta+\zeta^{-1}$ is $-\lambda$, so $\|v\|_1\le40\lambda^k$. Yet $|F_v(4)|=10\lambda^{-k}$ and $F_v$ vanishes on the other blocks. The two nonnegative vectors $v^+$ and $v^-$ display the loss of one power of precision in the next example.

\newpage
\section{A worked modulus and the scale at height \texorpdfstring{$T$}{T}}
For $q=11$, the element $2$ generates $G$ of order $10$. Its cyclic subgroups have orders $1,2,5,10$. Inverse pairs of generators and example choices are as follows.
\begin{center}
\begin{tabular}{cll}
\toprule
$d_H$ & Generators grouped by inversion & One class for exact data\\
\midrule
1  & $\{1\}$ & $1$\\
2  & $\{10\}$ & $10$\\
5  & $\{3,4\},\ \{5,9\}$ & $4$\\
10 & $\{2,6\},\ \{7,8\}$ & $2$\\
\bottomrule
\end{tabular}
\end{center}
The last two subgroups both have degree $4$. With $A=\{1,10,4,2\}$ the minimum spacing is of order $M^{-1}$. No fifth class improves this exponent, since at least one degree-four block still has a single observed pair. Adding \emph{both} $5$ and $8$ gives $\delta_A(M)=1$ for the enlarged set, for every $M\ge1$. More generally, uniform separation occurs precisely when all inverse pairs are represented. Its minimal sample size is $(|G|+|G[2]|)/2$, where $G[2]=\{a:a^2=1\}$.

For actual zeros at height $|\Im\rho|\le T$, the classical zero count in unit intervals gives $\sum_\chi m_\chi(\rho)\ll_q\log(T+3)$; see the count used in~\cite[\S2.3, \S2.7]{KR}. For fixed $q$ choose an integer cap $M\asymp_q\log(T+3)$ large enough. If each measured value in \eqref{eq:F} after normalization by $-h$ has error strictly below $\tfrac12(2M)^{-\kappa(A)}$, the nearest candidate in $\calM_M$ is unique and correct. Thus $\kappa(A)\log_2\log(T+3)+O_{q,A}(1)$ bits after the binary point suffice for either coordinate, with the additive constant accounting for normalization and complex absolute error. The matching upper bound on spacing describes worst-case integer vectors. Realization of those vectors by zeros is a separate number-theoretic problem.

\smallskip
\noindent\textbf{Connection and priority.} The readable Lean development~\cite{lean} passes from a one-sided remainder for weighted primes in progressions to zero-freeness. ProveIt's compressed prime-race packet~\cite{proveit} studies detecting zeros through weighted prime data. Here the data are individual progression residues and the target is the entire multiplicity vector. The sharp sample/precision curve above combines the classical grouped-zero formula and rational Fourier decomposition with a norm bound and a matching packing construction. We found no statement of this progression-by-progression curve in the sources below.

\begin{thebibliography}{9}\small
\bibitem{KR} H. Kisilevsky and M. O. Rubinstein, \emph{Chebotarev sets}, Acta Arith. 171 (2015), 97--124. \href{https://doi.org/10.4064/aa171-2-1}{doi:10.4064/aa171-2-1}.
\bibitem{PW} S. Perlis and G. L. Walker, \emph{Abelian group algebras of finite order}, Trans. Amer. Math. Soc. 68 (1950), 420--426. \href{https://doi.org/10.1090/S0002-9947-1950-0034758-3}{Publisher record}.
\bibitem{CO} W. R. Casper and B. Orozco, \emph{Reconstructing rational functions on finite Abelian groups with higher autocorrelations}, Thm.~2.8. \href{https://arxiv.org/abs/2503.21022}{arXiv:2503.21022}.
\bibitem{lean} L. Sernau et al., \href{https://github.com/lsernau/grh-lean/tree/72fead102f663d5c7a407964a13dc0cf9d977207}{\emph{grh-lean}}, commit \texttt{72fead10}, especially \texttt{OneSidedBound.lean}. \href{https://math.mit.edu/~drew/shortgrh/}{Independent verification record}.
\bibitem{proveit} ProveIt, \href{https://github.com/VladimirReshetnikov/ProveIt/tree/main/Analysis/GeneralizedRiemannHypothesis/Research/grh_relative_prime_races_2026-11-25}{\emph{Relative zeta zeros from compressed prime races}}, 25 November 2026.
\end{thebibliography}
\end{document}
